% Рисунок из тетради: сходимость в D(Omega) — носители в одном компакте K.
\begin{figure}[htbp]
\centering
\begin{tikzpicture}
\begin{axis}[
  width=10cm, height=4.4cm,
  axis lines=middle,
  xmin=-2.6, xmax=2.6, ymin=-0.55, ymax=2.3,
  xtick=\empty, ytick=\empty,
  samples=120, clip=false,
]
  % общий компакт K
  \draw[very thick, gray!70] (axis cs:-1.5,-0.22) -- (axis cs:1.5,-0.22);
  \draw[very thick, gray!70] (axis cs:-1.5,-0.1) -- (axis cs:-1.5,-0.34);
  \draw[very thick, gray!70] (axis cs: 1.5,-0.1) -- (axis cs: 1.5,-0.34);
  \node[gray, anchor=north, font=\footnotesize] at (axis cs:0,-0.4) {$K$};

  % предельная функция и приближения — все с носителем внутри K
  \addplot[very thick, blue!60!black, domain=-1.2:1.2] {1.7*bump(x,1.2)/bump(0,1.2)};
  \addplot[thick, red!70!black, domain=-0.9:0.9]  {1.25*bump(x-0.15,0.9)/bump(0,0.9)};
  \addplot[thick, orange!80!black, domain=-1.1:1.1] {1.95*bump(x+0.1,1.1)/bump(0,1.1)};

  \node[blue!60!black, anchor=west, font=\footnotesize] at (axis cs:1.25,1.1) {$\varphi$};
  \node[red!70!black, anchor=west, font=\footnotesize] at (axis cs:1.6,0.55) {$\varphi_k$};
\end{axis}
\end{tikzpicture}
\caption{Сходимость $\varphi_k \xrightarrow{\mathcal{D}} \varphi$: все носители лежат
в \emph{одном} компакте $K \subset \Omega$, и каждая производная $D^\alpha \varphi_k$
сходится к $D^\alpha\varphi$ равномерно. Без общего компакта носители могли бы
«уползать» к границе $\Omega$.}
\label{fig:test-conv}
\end{figure}
